\chapter{iCub humanoid robot}


\chapter{Introduction to rigid multibody dynamics}
Intro here

\section{Notation overview}
The following mathematical notation is used throughout the thesis.
\begin{itemize}
    \item The set of real numbers is $\mathbb{R}$. If $u$ and $v$ are two $n$-dimensional column vectors (e.g. $u, v \inrn{n}$), then their inner product is $u^Tv$, where $'T'$ is the transpose operator.
    \item The identity matrix of dimension $n$ is denoted as $I_{n} \inrn{n\times n}$. The zero column vector of dimension $n$ is denoted $0_n \inrn{n}$, while the zero matrix of dimension $n\times m$ is denoted $0_{n\times m} \inrn{n\times m}$.
    \item The set $SO(3)$, also called 3D rotation group, is a set of all rotations about the origin of space $\mathbb{R}^3$ under the operation of composition. In other words, it is the set of $\mathbb{R}^{3\times 3}$ orthogonal matrices with determinant equal to one:

    \begin{equation}\label{def_SO3_group}
        SO(3) \defeq \big\{ R \inrn{3\times 3} \;|\; R^TR = I_3,\; \text{det}(R) = 1 \big\}.
    \end{equation}

    \item The set $\mathfrak{so}(3)$ is the set of $3\times 3$ skew-symmetric matrices:
    \begin{equation}\label{def_so3_group}
        \mathfrak{so}(3) \defeq \big\{ S \inrn{3\times 3} \;|\; S^T = -S \big\}
    \end{equation}

    \item The set $SE(3)$ is defined as
    \begin{equation}\label{def_SE3_group}
        SE(3) \defeq \Big\{ 
        \begin{bmatrix}
            R & p\\ 0_{1\times3} & 1
        \end{bmatrix} 
        \inrn{4\times 4} \;|\; R\in SO(3), p\inrn{3}
        \Big\}
    \end{equation}.

    \item The set $\mathfrak{se}(3)$ is defined as
    \begin{equation}\label{def_se3_group}
        \mathfrak{se}(3) \defeq \Big\{ 
        \begin{bmatrix}
            \Omega & v\\ 0_{1\times3} & 0
        \end{bmatrix} 
        \inrn{4\times 4} \;|\; \Omega\in \mathfrak{so}(3), v\inrn{3}
        \Big\}
    \end{equation}. 

    \item Given the vector $w = (x,y,z) \inrn{3}$, we define $\sskew{w}$ as the $3\times3$ skew-symmetric matrix
    \begin{equation}\label{def_skew_symmetric_matrix}
        \sskew{w} = 
        \begin{bmatrix}
            x\\y\\z
        \end{bmatrix}^\wedge
     \defeq
    \begin{bmatrix}
        0 & -z & y\\
        z & 0 & -x\\
        -y & x & 0
    \end{bmatrix}
    \in \mathfrak{so}(3)
    \end{equation}

    \item Given a vector $v = (v; w) \inrn{6}$ ($v, \omega \inrn{3}$), we define
    \begin{equation}
        \sskew{\vv} =
        \begin{bmatrix}
            v\\ \omega
        \end{bmatrix}^\wedge \defeq
        \begin{bmatrix}
            \sskew{\omega} & v\\
            0_{1\times3} & 0
        \end{bmatrix}
        \in \mathfrak{se}(3)
    \end{equation}.

    \item Given a vector $v = (v; w) \inrn{6}$ ($v, \omega \inrn{3}$), we define
    \begin{equation}
        \vv\times \defeq
        \begin{bmatrix}
            \sskew{\omega} & \sskew{v}\\
            0_{1\times3} & \sskew{\omega}
        \end{bmatrix}
        \inrn{6\times6}
    \end{equation}

    \item Given a vector $v = (v; w) \inrn{6}$ ($v, \omega \inrn{3}$), we define
    \begin{equation}
        \vv\overline{\times}^{*} \defeq
        \begin{bmatrix}
            \sskew{\omega} & 0_{3\times 3}\\
            \sskew{v} & \sskew{\omega}
        \end{bmatrix}
        \inrn{6\times6}
    \end{equation}

    \item For a matrix $A \inrn{n\times m}$ and $B\inrn{p\times q}$, the Kronecker product is $A\otimes B \inrn{np\times mq}$
\end{itemize}

\section{Coordinate frames and transformations}
In rigid-body dynamics, coordinate frames serve as fundamental reference systems for describing the position and orientation of objects in three-dimensional space. The correlation between two coordinate frames is generally defined using homogeneous transformations, which combine rotation and translation into a single $4 \times 4$ matrix. This method provides a simple and unified structure for performing coordinate transformations efficiently using matrix multiplication. The rotational component of a homogeneous transformation belongs to the special orthogonal group $\text{SO}(3)$, ensuring that rigid body rotations preserve distances and angles. The translational component, represented as a column vector, captures the displacement between the two frames. 

A frame $A$ can then be expressed as $A = (o_A, [A])$, where $o_A$ is the origin of the frame and $[A]$ is its
orientation. In this way, we can obtain the coordinate vector $\leftsuperscript{A}{p}{} \in \mathbb{R}^3$ expressed in the frame orientation $[A]$ as follows:

\[
\leftsuperscript{A}{p}{} = \begin{bmatrix}
    \vec{r_{o_A,p}} \cdot \vec{x_A} \\ 
    \vec{r_{o_A,p}} \cdot \vec{y_A} \\
    \vec{r_{o_A,p}} \cdot \vec{z_A}.
\end{bmatrix}
\]

Given two frames A and B, we can define the coordinate transformation from B to A as $\leftsuperscript{A}{R}{B} \in \text{SO}(3)$, which only depends on the relative orientation
between the orientation frames [A] and [B], and not on the positions of their origins. If we would like to describe both the position and orientation of frame $B$ w.r.t. frame $A$ in a compact form, we can use the $4 \times 4$ homogeneous transformation matrix mentioned earlier, also known as transform:

\[
\leftsuperscript{A}{H}{B} = \begin{bmatrix}
    \leftsuperscript{A}{R}{B} & \leftsuperscript{A}{o}{B} \\
    0_{1\times 3} & 1
\end{bmatrix} \in \text{SE}(3).
\]

Given a point $p$, the homogeneous transformation matrix $\leftsuperscript{A}{H}{B}$ can be also
used to map the coordinate vector $\leftsuperscript{A}{p}{}$ to $\leftsuperscript{B}{p}{}$:

\[
\leftsuperscript{A}{\tilde{p}}{} = \leftsuperscript{A}{H}{B} \leftsuperscript{B}{\tilde{p}}{}, ¯
\]
where $\leftsuperscript{A}{\tilde{p}}{}$ is the homogeneous representation of the coordinate vector $\leftsuperscript{A}{p}{}$. That is, $\leftsuperscript{A}{\tilde{p}}{} = (\leftsuperscript{A}{p}{}; 1)$ (where ; indicates row concatenation). In matrix form, we get $\leftsuperscript{A}{p}{} = \leftsuperscript{A}{R}{B} \leftsuperscript{B}{p}{} + \leftsuperscript{A}{o}{B}$.
In can also be shown that, given a transform $\leftsuperscript{A}{H}{B}$, the inverse transformation can be expressed as:

\begin{equation}
    \leftsuperscript{B}{H}{A} = \leftsuperscript{A}{H}{B}^{-1} = \begin{bmatrix}
    \leftsuperscript{B}{R}{A} & \leftsuperscript{B}{o}{A} \\
    0_{1\times 3} & 1
\end{bmatrix} = \begin{bmatrix}
    \leftsuperscript{A}{R}{B} & -\leftsuperscript{B}{R}{A}\leftsuperscript{A}{o}{B} \\
    0_{1\times 3} & 1
\end{bmatrix}
\end{equation}


\section{Velocities}
Given a point $p$ and a frame A, we can define:

\begin{equation}
    \leftsuperscript{A}{\dot{p}}{} \defeq \ddt{\leftsuperscript{A}{p}{}}.
\end{equation}

In particular, when $p$ is the origin of a frame, e.g., $p = o_B$, we have:

\[
\leftsuperscript{A}{\dot{o}}{B} = \ddt{\leftsuperscript{A}{o}{B}}.
\]

In the same way, we can define

\begin{equation}
    \leftsuperscript{A}{\dot{R}}{B} \defeq \ddt{\leftsuperscript{A}{R}{B}},
\end{equation}
and
\begin{equation}
    \leftsuperscript{A}{\dot{H}}{B} \defeq \ddt{\leftsuperscript{A}{H}{B}} = 
    \begin{bmatrix}
    \leftsuperscript{A}{\dot{R}}{B} & \leftsuperscript{A}{\dot{o}}{B} \\
    0_{1\times 3} & 0
    \end{bmatrix}.
\end{equation}

A more compact representation of $\leftsuperscript{A}{\dot{H}}{B}$ can be obtained by multiplying it with the inverse of $\leftsuperscript{A}{H}{B}$ on the left or on the right. In both cases, the result is an element of the $\mathfrak{se}(3)$ that will be called a \textit{6D velocity}.

\subsection{Left-trivialised velocity}

By multiplying it on the left:

\begin{equation} \label{def_left_trivialised_homogeneous}
    {}^A H_B^{-1} \dot{{}^A H_B} =
    \begin{bmatrix}
    {}^A R_B^T & -{}^A R_B^T {}^A o_B \\
    0_{1\times3} & 1
    \end{bmatrix}
    \begin{bmatrix}
    {}^A \dot{R}_B & {}^A \dot{o}_B \\
    0_{1\times3} & 0
    \end{bmatrix}
    =
    \begin{bmatrix}
    {}^A R_B^T {}^A \dot{R}_B & {}^A R_B^T {}^A \dot{o}_B \\
    0_{1\times3} & 0
    \end{bmatrix} \in \mathfrak{se}(3).
\end{equation}

Note that $\leftsuperscript{A}{R}{B}^{T} \leftsuperscript{A}{\dot{R}}{B}$ appearing on the right hand side of \ref{def_left_trivialised_homogeneous} is skew symmetric. We can define $\leftsuperscript{B}{v}{A,B}$ and $\leftsuperscript{B}{\omega}{A,B} \in \mathbb{R}^3$ so that:

\begin{equation}\label{def_linear_velocity}
    \leftsuperscript{B}{\varv}{A,B} \defeq \leftsuperscript{A}{R}{B}^{T} \leftsuperscript{A}{\dot{o}}{B},
\end{equation}
\begin{equation}\label{def_angular_velocity}
    \leftsuperscript{B}{\omega}{A,B}^{\wedge} \defeq \leftsuperscript{A}{R}{B}^{T} \leftsuperscript{A}{\dot{R}}{B}.
\end{equation}

Equations \ref{def_linear_velocity} and \ref{def_angular_velocity} respectively define the linear and angular velocity of frame B with respect to frame A.
We can merge these together by defining the \textit{left-trivialised} velocity of frame B with respect to frame A as:

\begin{equation}\label{def_left_trivialised_velocity}
    \leftsuperscript{B}{v}{A,B} = \begin{bmatrix}
        \leftsuperscript{B}{\varv}{A,B} \\
        \leftsuperscript{B}{\omega}{A,B}
    \end{bmatrix}
    \in \mathbb{R}^6.
\end{equation}

Then, by construction from equation \ref{def_left_trivialised_homogeneous}, we have:

\begin{equation} \label{left_trivialised_construction}
    \leftsuperscript{B}{v}{A,B}^{\wedge} = \leftsuperscript{A}{H}{B}^{-1} \dot{\leftsuperscript{A}{H}{B}}.
\end{equation}

\subsection{Right-trivialised velocity}

Similarly to what was done for equation \ref{def_left_trivialised_homogeneous}, we can multiply by $\leftsuperscript{A}{H}{B}^{-1}$ on the right side of $\leftsuperscript{A}{\dot{H}}{B}$ instead of the left side, to obtain:

\begin{equation}\label{def_right_trivialised_homogeneous}
    \leftsuperscript{A}{\dot{H}}{B} \leftsuperscript{A}{H}{B}^{-1} = 
    \begin{bmatrix}
        \leftsuperscript{A}{\dot{R}}{B}^{T} & \leftsuperscript{A}{\dot{o}}{B} \\
        0_{1\times 3} & 0
    \end{bmatrix}
    \begin{bmatrix}
        \leftsuperscript{A}{R}{B}^{T} & -\leftsuperscript{A}{R}{B}^T \leftsuperscript{A}{o}{B} \\
        0_{1\times 3} & 1
    \end{bmatrix}
    = 
    \begin{bmatrix}
        \leftsuperscript{A}{\dot{R}}{B}\leftsuperscript{A}{R}{B}^{T} & \leftsuperscript{A}{\dot{o}}{B} - \leftsuperscript{A}{\dot{R}}{B}\leftsuperscript{A}{R}{B}^T\leftsuperscript{A}{o}{B} \\
        0_{1\times 3} & 0
    \end{bmatrix}
\end{equation}

Again, defining $\leftsuperscript{A}{\varv}{A,B}$ and $\leftsuperscript{A}{\omega}{A,B} \in \mathbb{R}^3$ as 

\begin{equation}
    \leftsuperscript{A}{\varv}{A,B} = \leftsuperscript{A}{\dot{o}}{B} - \leftsuperscript{A}{\dot{R}}{B}\leftsuperscript{A}{R}{B}^{T} \leftsuperscript{A}{\dot{o}}{B},
\end{equation}
\begin{equation}
    \leftsuperscript{A}{\omega}{A,B}^{\wedge} = \leftsuperscript{A}{\dot{R}}{B}\leftsuperscript{A}{R}{B}^{T},
\end{equation}

the \textit{right-trivialised velocity} of B with respect to A is defined as

\begin{equation} \label{def_right_trivialised_velocity}
    \leftsuperscript{A}{v}{A,B} = 
    \begin{bmatrix}
        \leftsuperscript{A}{\varv}{A,B} \\
        \leftsuperscript{A}{\omega}{A,B}
    \end{bmatrix}
    \in \mathbb{R}^6
\end{equation}

By construction,
\begin{equation} \label{right_trivialised_construction}
    \leftsuperscript{A}{v}{A,B}^{\wedge} = \leftsuperscript{A}{H}{B} \leftsuperscript{B}{v}{A,B}^{\wedge}\leftsuperscript{A}{H}{B}^{-1}
\end{equation}

\subsection{6D velocities in different frames}

By combining the definitions for the left and right-trivialised velocities, respectively described by equations \ref{left_trivialised_construction} and \ref{right_trivialised_construction}, we obtain the mapping between these velocities:

\begin{equation}
    \leftsuperscript{A}{v}{A,B}^{\wedge} = \leftsuperscript{A}{H}{B}\leftsuperscript{B}{v}{A,B}^{\wedge}\leftsuperscript{A}{H}{B}^{-1},
\end{equation}
which clearly is a linear mapping.
We can write the 6D velocity in linear form:

\begin{equation}\label{6d_velocity_linear_mapping}
    \leftsuperscript{A}{v}{A,B} = \leftsuperscript{A}{X}{B}\leftsuperscript{A}{v}{A,B},
\end{equation}

with

\begin{equation}\label{def_adjoint_matrix}
    \leftsuperscript{A}{X}{A,B} = \leftsuperscript{A}{X}{B} \Big(\leftsuperscript{A}{H}{B} \Big) = 
    \begin{bmatrix}
        \leftsuperscript{A}{R}{B} & \leftsuperscript{A}{o}{B}^{\wedge} \\
        0_{3 \times 3} & \leftsuperscript{A}{R}{B},
    \end{bmatrix} \in \mathbb{R}^{6 \times 6}
\end{equation}
which is typically called the \textit{adjoint} matrix.

Recalling that $\leftsuperscript{A}{o}{B} = -\leftsuperscript{A}{R}{B}\leftsuperscript{B}{o}{A}$, it can be shown that the inverse velocity transformation is $\leftsuperscript{B}{X}{A} = \leftsuperscript{A}{X}{B}^{-1}$.

\subsection{Mixed velocity representation}
One can observe that in neither representations of the velocity described earlier the linear part of the velocity $\varv$ is the derivative of the position vector of the origin of frame B ($\leftsuperscript{A}{\dot{o}}{B}$). However, in most situations, it is convenient to adopt this notation. We can do this by introducing the frame $B[A] \defeq \Big( o_B, [A] \Big)$ as the frame with the same origin of B and the same orientation of A. We have then:

\begin{equation}\label{def_mixed_velocity_frame}
    \leftsuperscript{B[A]}{H}{B} = 
    \begin{bmatrix}
        \leftsuperscript{A}{R}{B} & 0_{3\times 1}\\
        0_{3\times 1} & 0
    \end{bmatrix},
\end{equation}

and by expressing the velocity $v_{A,B}$ in $B[A]$, we get:

\begin{equation}\label{def_mixed_velocity}
    \leftsuperscript{B[A]}{v}{A,B} = \leftsuperscript{B[A]}{X}{B}\leftsuperscript{B}{v}{A,B} = 
    \begin{bmatrix}
        \leftsuperscript{A}{R}{B} & 0\\
        0 & \leftsuperscript{A}{R}{B}
    \end{bmatrix}
    \begin{bmatrix}
        \leftsuperscript{B}{R}{A}\leftsuperscript{A}{\dot{o}}{B}\\
        \leftsuperscript{B}{\omega}{A,B}
    \end{bmatrix} = 
    \begin{bmatrix}
        \leftsuperscript{A}{\dot{o}}{B}\\
        \leftsuperscript{A}{\omega}{A,B}
    \end{bmatrix}
\end{equation}

\subsection{The cross product on $\mathbb{R}^6$}
Equation \ref{left_trivialised_construction} can be rewritten as
\begin{equation}
    \leftsuperscript{A}{\dot{H}}{B} = \leftsuperscript{A}{H}{B}\leftsuperscript{B}{v}{A,B}^{\wedge}
\end{equation}
Differentiating equation \ref{def_adjoint_matrix} with respect to time, we have:

\begin{equation} \label{def_adjoint_matrix_time_derivative}
    \leftsuperscript{A}{\dot{X}}{B} =
    \begin{bmatrix}
        \leftsuperscript{A}{\dot{R}}{B} & \leftsuperscript{A}{\dot{o}}{B}^{\wedge}\leftsuperscript{A}{R}{B} + \wedged{\leftsuperscript{A}{O}{B}}\leftsuperscript{A}{\dot{R}}{B}\\
        0_{3\times 3} & \leftsuperscript{A}{\dot{R}}{B}
    \end{bmatrix} = \leftsuperscript{A}{X}{B} \leftsuperscript{B}{v}{A,B}^{\times},
\end{equation}
with $\leftsuperscript{B}{v}{A,B}^{\times}$ defined as
\begin{equation}\label{def_cross_product_r6}
    \leftsuperscript{B}{v}{A,B}^{\times} \defeq
    \begin{bmatrix}
    \wedged{\leftsuperscript{B}{\omega}{A,B}} & \wedged{\leftsuperscript{B}{\varv}{A,B}}\\
    0_{3\times 3} & \wedged{\leftsuperscript{B}{\omega}{A,B}}
    \end{bmatrix}.
\end{equation}
Equation \ref{def_cross_product_r6} is the matrix representation of the cross product in $\mathbb{R}^6$.


\section{Accelerations}
The acceleration of a frame B relative to another frame A expressed in a generic frame C is obtaind by taking the tme-derivative of the associated velocity:

\begin{equation}\label{def_frame_acceleration}
    \leftsuperscript{C}{\dot{v}}{A,B} = \ddt{\leftsuperscript{C}{v}{A,B}}.
\end{equation}

Extracting the left-trivialised velocity and the time-derivative of the adjoint matrix (\ref{def_adjoint_matrix_time_derivative}), we obtain:

\begin{equation}
    \leftsuperscript{C}{\dot{v}}{A,B} = \ddt{\leftsuperscript{C}{X}{B}\leftsuperscript{B}{v}{A,B}} = \leftsuperscript{C}{X}{B}\leftsuperscript{B}{\dot{v}}{A,B} + \leftsuperscript{C}{\dot{X}}{B}\leftsuperscript{B}{v}{A,B} = \leftsuperscript{C}{X}{B} \Big( \leftsuperscript{B}{\dot{v}}{A,B} + \leftsuperscript{B}{v}{C,B}^{\times}\leftsuperscript{B}{v}{A,B} \Big).
\end{equation}

It can be noted that the simpler form $\leftsuperscript{C}{\dot{v}}{A,B} = \leftsuperscript{C}{X}{B}\leftsuperscript{B}{\dot{v}}{A,B}$ does not hold. If, however, either $C = A$ or $C = B$, the cross-product term is zero and the equality holds.
Under these conditions, the accelerations are

\begin{equation}
    \leftsuperscript{C}{a}{A,B} = \leftsuperscript{C}{X}{A}\leftsuperscript{A}{\dot{v}}{A,B} = \leftsuperscript{C}{X}{B}\leftsuperscript{B}{\dot{v}}{A,B}.
\end{equation}

The acceleration $\leftsuperscript{C}{\dot{v}}{A,B}$ is called \textit{apparent acceleration}, while $\leftsuperscript{C}{a}{A,B}$ is called \textit{intrinsic acceleration}.
We can also call \textit{proper acceleration} as the intrinsic acceleration without the effects of gravitational accelerations, defined as follows:

\begin{equation} \label{def_proper_acceleration}
    \leftsuperscript{C}{\overline{a}}{A,B} = \leftsuperscript{C}{a}{A,B} - \leftsuperscript{C}{X}{A}\begin{bmatrix}
        \leftsuperscript{A}{R}{W}\leftsuperscript{W}{g}{}\\
        0_{3 \times 1}
    \end{bmatrix}.
\end{equation}
Here the definition of gravity is $\leftsuperscript{W}{g}{} = (0, 0, -g) \in \mathbb{R}^3$.

\section{Forces \& Torques}
The force acting on a rigid body can be expressed as a stacking of a linear component $\mathcal{f} \in \mathbb{R}^3$ and an angular component $\mathcal{\tau} \in \mathbb{R}^3$. We can then define the force $f \in \mathbb{R}^6$ with respect to a frame B as

\begin{equation} \label{def_force_6d}
    {}_B\text{f} \defeq
    \begin{bmatrix}
        \leftsuperscript{B}{f}{}\\
        \leftsuperscript{B}{\mathcal{\tau}}{}
    \end{bmatrix}.
\end{equation}

Similarly to the work done for the 6D velocity, we can define a linear map to change the coordinates of the 6D force from a frame B to another frame A. The transformation is ${}_AX^{B}$, and written as

\begin{equation} \label{force_change_frame}
    {}_A\text{f} = {}_AX^{B} {}_B\text{f}
\end{equation}

The mapping ${}_AX^{B}$, strictly related to the transformation of 6D velocities between the same frames, can be defined as 

\begin{equation} \label{mapping_definition_transpose}
{}_AX^B = \leftsuperscript{B}{X}{A}^T =
\begin{bmatrix}
    \leftsuperscript{A}{R}{B} & 0_{3\times 3}\\
    \leftsuperscript{A}{\hat{o}}{B}\leftsuperscript{A}{R}{B} & \leftsuperscript{A}{R}{B}
\end{bmatrix}
\in \mathbb{R}^{6\times 6}
\end{equation}

It is important to realize that \ref{mapping_definition_transpose} is such to make the following identity hold:

\begin{equation}
    \langle {}_B\text{f}, \leftsuperscript{B}{v}{A,B}\rangle = \langle {}_A\text{f}, \leftsuperscript{A}{v}{A,B} \rangle,
\end{equation}

where f can be interpreted as a 6D force applied to a rigid body to which the moving frame B is rigidly attached and A as the absolute inertial frame.\\



\section{Rigid body Lagrangian dynamics}

Given a body B, in order to simplify the notation, we denote the coordinates of a point $p_i$ belonging to the body and expressed in B as $r = \leftsuperscript{B}{p}{i}$. 
Some useful concepts that help us to define the lagrangian dynamics of the rigid body are:
\begin{itemize}
    \item The \textit{total mass} of the rigid body is calculated by introducing the function $\rho: \mathbb{R}^3 \longrightarrow \mathbb{R}^+$, called density function, that maps each point of the body to its density, and then integrating over the volume occupied by the body:

    \[
    m = \iiint_{\mathbb{R}^3} \rho(r)dr \in R
    \]

    \item The \textit{centre of mass} (CoM) of the rigid body is the average point of its density:

    \begin{equation}\label{centre_of_mass_rigid_body}
        c = \leftsuperscript{B}{p}{CoM} = \frac{\iiint_{\mathbb{R}^3} r\rho(r)dr }{\iiint_{\mathbb{R}^3} \rho(r)dr} = \frac{1}{m} \iiint_{\mathbb{R}^3} r\rho(r)dr \inrn{3}
    \end{equation}

    \item The \textit{inertia tensor} of the rigid body, describing all moments of inertia of a body rotating around a specific axis, can be computed as:

    \begin{equation}\label{inertia_tensor_rigid_body}
        I = - \iiint_{\mathbb{R}^3} \rho(r)(r^{\wedge})^2 dr \inrn{3\times3}
    \end{equation}
\end{itemize}


The Lagrangian provides a powerful method for deriving the equations of motion of a system without directly using Newton’s laws. In fact, it uses energy definitions rather than forces explicitly.
To obtain the equations of motion for the rigid body through a Lagrangian approach, it is useful to also define a few functions describing the energy of the body. If we consider a rigid body B with pose $H \in SE(3)$ and velocity $\dot{H}$, the Lagrangian function of B is

\begin{equation}\label{def_lagrangian_rb}
    L(H,\dot{H}) = K(H,\dot{H}) - U(H),
\end{equation}
\begin{equation}\label{def_kinetic_energy}
    K(H, \dot{H})) = \frac{1}{2} \iiint_{\mathbb{R}^3} \rho(r) \Big| \dot{R}r + \dot{o} \Big|^2 dr,
\end{equation}
\begin{equation} \label{def_potential_energy}
    U(H) = \iiint_{\mathbb{R}^3} \rho(r)g^T (Rr + o) dr.
\end{equation}
The change of variables $p = Rr + o, \dot{p} = \dot{R}r + \dot{o}$ has been made to highlight the dependency of the Lagrangian function on the state of the rigid body $H$ and $\dot{H}$.

A quick note describes two main differences between the Lagrangian of a point particle $L(p,\dot{p}) = K(p, \dot{p}) - U(p)$ and of a rigid body \ref{def_lagrangian_rb}. First, the configuration of a point particle $p \inrn{3}$ is simply an element of a vector space, while the state of the rigid
body is an element of a matrix Lie group 
$
H = \begin{bmatrix}
    R & o\\
    0_{1\times3} & 1
\end{bmatrix} \in SE(3),
$
which we previously called \textit{pose} of a rigid body.

Secondly, the parameter describing the inertia of a point mass particle is a single lumped single parameter, the mass of the point $m$. For a rigid body, the inertia information is contained in the density function $\rho(\cdot)$.
The first difference is tackled using the left-trivialized velocity of a rigid body, while the second difference by introducing the inertial parameters, a set of 10 parameters fully describing all the inertial properties of a rigid body.
To transform $\dot{H}$ in a 6D vector we opt to use the left-trivialized velocity because we can directly apply the Euler-Poincarè equations.
If we refer to the left-trivialised velocity $\leftsuperscript{B}{v}{A,B}$ as $v$:

\begin{equation} \label{left_trivialised_velocity_lagrangian}
    \vv = \begin{bmatrix}
        v\\w
    \end{bmatrix} = 
    \begin{bmatrix}
        R^T\dot{o}\\
        (R^T\dot{R})^{\wedge},
    \end{bmatrix},
    \quad
    \dot{H}=H\vv^{\wedge}=
    \begin{bmatrix}
        R\omega^{\wedge} & Rv\\
        0_{3\times1} & 0
    \end{bmatrix},
\end{equation}
we can express the Lagrangian \ref{def_lagrangian_rb} in function of the left-trivialized
velocity:


\begin{equation} \label{left_trivialised_lagrangian}
    l(H, \vv) = L(H, H\vv^{\wedge}) = k(H,\vv) - U(H),
\end{equation}
\begin{equation}\label{left_trivialised_kinetic_energy}
    \begin{split}
        k(H, \vv) = \frac{1}{2}\iiint_{\mathbb{R}^3} \rho(r) \Big| R\omega^{\wedge}r + Rv \Big|^2 dr =\\ =\frac{1}{2}\iiint_{\mathbb{R}^3} \rho(r) \Big| R(\omega^{\wedge}r + v  ) \Big|^2 dr =\\= \frac{1}{2}\iiint_{\mathbb{R}^3} \rho(r) \Big| \omega^{\wedge}r + v \Big|^2 dr
    \end{split}
\end{equation}

The \textit{left trivialised Lagrangian} of a rigid body can be written as

\begin{equation}\label{left_trivialised_lagrangian2}
    l(H,\vv) = K(\vv) - U(H),
\end{equation}
\begin{equation}\label{left_trivialised_kinetic_energy2}
    k(\vv) = \frac{1}{2} \vv^T \mathbb{M} \vv,
\end{equation}
\begin{equation}\label{left_trivialised_potential_energy2}
    U(H) = \begin{bmatrix}
        g^T & 0
    \end{bmatrix} H \begin{bmatrix}
        mc\\
        m
    \end{bmatrix},
\end{equation}
where $m\in \mathbb{R}$ is the total mass of the body, $c \inrn{3}$ is the centre of mass of the body and $I \inrn{3\times 3}$ is the 3D inertia matrix of the body as described before. Then $\mathbb{M}\inrn{6\times6}$ is the inertia matrix:

\begin{equation}\label{def_inertia_matrix}
    \mathbb{M} = \begin{bmatrix}
        m{1}_3 & -(mc)^{\wedge}\\
        (mc)^{\wedge} & I
    \end{bmatrix}.
\end{equation}
This simplified version of the left-trivialised Lagrangian allows to remove the dependency on the density function, but depends on a fixed number of \textit{inertial parameters}.

The system's positions and velocities are now connected by a relationship between $\leftsuperscript{A}{H}{B} \in SE(3)$ and $\vv_{A,B}^{\wedge}$, which allows the usage of the Euler-Poincarè equation to obtain the system's dynamics.
Let us define a quantity called the \textit{action} S, defined as

\begin{equation} \label{def_action}
    S[H(\cdot)] = \int_0^T L(H(t), \dot{H}(t)) dt,
\end{equation}
where the rigid body Lagrangian $L(H, \dot{H})$ has been defined in (\ref{def_lagrangian_rb}).
The trajectory $H(\cdot)$ of a rigid in a uniform gravitational field is the one that minimises the action. 

The resulting equations of motion are given by:

\begin{subnumcases}
    \dot{H} = H\vv^\wedge \label{euler_poincare_a} \\
    \mathbb{M}\dot{\vv} + \vv\overline{\times}^{*} \mathbb{M}\vv = \mathbb{M}\begin{bmatrix}
        R^T g\\
        0_{3\times 1}
    \end{bmatrix}
    + {}_B\text{f}^{\text{ext}} \label{euler_poincare_b}
\end{subnumcases}
    
where $\mathbb{M}$ is the 6D inertia matrix defined in \ref{def_inertia_matrix}, $\vv\overline{\times}^{*}$ is defined in \ref{def_cross_product_r6} and ${}_B\text{f}^{\text{ext}}$ is the sum of the external non-gravitational force-torques acting in the body frame $B$, which make the equations of motion more generic.

\subsection{Newton-Euler equations in different representations}
By applying the velocity and accelerations transformations previously defined, it is possible to transform the Newton-Euler equations to be expressed with respect to different velocity/acceleration representations.

\subsubsection{Left-trivialised}
Using the left-trivialised form of the acceleration, equation \ref{euler_poincare_b} becomes:

\begin{equation} \label{euler_poincare_left_trivialised}
    {}_B\mathbb{M}_B \leftsuperscript{B}{\dot{\vv}}{A,B} + \leftsuperscript{B}{v}{A,B}\overline{\times}^{*} {}_B\mathbb{M}_B \leftsuperscript{B}{\vv}{A,B} = {}_B\mathbb{M}_B
    \begin{bmatrix}
        \leftsuperscript{A}{R}{B}^T \leftsuperscript{A}{g}{}\\
        0_{3\times1}
    \end{bmatrix}
    + {}_B\text{f}^{\textbf{ext}}
\end{equation}

\subsubsection{Right-trivialised}

\begin{equation} \label{euler_poincare_right_trivialised}
    {}_A\mathbb{M}_B \leftsuperscript{A}{\dot{\vv}}{A,B} + \leftsuperscript{A}{v}{A,B}\overline{\times}^{*} {}_A\mathbb{M}_B \leftsuperscript{A}{\vv}{A,B} = {}_A\mathbb{M}_B
    \begin{bmatrix}
        \leftsuperscript{A}{g}{}\\
        0_{3\times1}
    \end{bmatrix}
    + {}_A\text{f}^{\textbf{ext}}
\end{equation}

Despite this equation may appear similar to (\ref{euler_poincare_left_trivialised}), the main difference lies in the fact that ${}_B\mathbb{M}_B$ s a fixed quantity, while ${}_A\mathbb{M}_B$ is time-varying and depends on the position of the rigid body.

\subsubsection{Mixed}
The dynamics in mixed acceleration is

\begin{equation}\label{euler_poincare_mixed}
    {}_{B[A]}\mathbb{M}_B \leftsuperscript{B[A]}{\dot{\vv}}{A,B} + 
    \begin{bmatrix}
        0_{3\times 1}\\
        \leftsuperscript{A}{\omega}{A,B}
    \end{bmatrix}
    \overline{\times}^{*} {}_{B[A]}\mathbb{M}_B 
    \begin{bmatrix}
        0_{3\times1}\\
        \leftsuperscript{A}{\omega}{A,B}
    \end{bmatrix}
    = {}_{B[A]}\mathbb{M}_B
    \begin{bmatrix}
        \leftsuperscript{A}{g}{}\\
        0_{3\times1}
    \end{bmatrix}
    + {}_{B[A]}\text{f}^{\textbf{ext}}
\end{equation}


\section{Multi-body Lagrangian dynamics}
\subsection{Introduction to multi-body systems}
A multi-body system consists of multiple interconnected rigid bodies that undergo motion while interacting through forces and constraints. These systems are widely used to model mechanical, robotic, biomechanical and other applications, where understanding the motion and interactions between bodies is essential for the design and functionality. The study of multi-body dynamics provides mathematical a strong framework to describe the kinematics and dynamics of these systems, considering factors such as joint constraints, external forces, and system topology. By applying principles from classical mechanics, such as Newton-Euler and Lagrangian formulations, multi-body dynamics enables accurate simulation and optimization of complex mechanical systems.
To be more specific, a multi-body system is composed of different components.


\subsubsection{Joints}\label{sssec_joints}
A multi-body system is composed of $n_j$ joints. A \textit{joint} is a connection that constrains the relative motion between two rigid bodies. The number of degrees of freedom (dofs) of a joint is the number of relative degrees of freedom between the two bodies that are left unconstrained by the joint, which typically range from 0 to 6. In this chapter we generally consider joints which have only 1 degree of freedom, so $n_j = n_{dof} = n$.
For one degree of freedom joints, the transform between the two connected bodies $B$ and $D$ is fully determined by the function mapping the joint position $\theta \in \mathbb{R}$ to the relative pose of the two links:
\[
\leftsuperscript{B}{H}{D}(\theta) :\mathbb{R} \mapsto SE(3).
\]
The most common types of joints with 1 dof are called \textit{revolute} and \textit{prismatic} joints.
The relative velocity and acceleration of D with respect to B is related to the time derivatives of the joint configuration $\theta$:
\begin{equation}
 \iddt{\leftsuperscript{B}{H}{D}(\theta)} = \frac{d \Big(\leftsuperscript{B}{H}{D}(\theta) \Big)}{d \theta}\dot{\theta}
\end{equation}
It is useful to express the relative velocity as a 6-dimensional vector. If we call $X$ the placeholder that selects any velocity representations introduced in chapter(FILL), the relative velocity between links $B$ and $D$ is:

\begin{equation}\label{relative_velocity_joint_motion_subspace}
    \leftsuperscript{X}{\vv}{B,D} = \leftsuperscript{X}{S}{B,D}(\theta)\dot{\theta}
\end{equation}
where we introduced the \textit{joint motion subspace vector} $\leftsuperscript{X}{S}{B,D}(\theta) \inrn{6}$.
In different velocity representations:

\[
\leftsuperscript{D}{S}{B,D}(\theta) = \Big[\leftsuperscript{D}{H}{B}(\theta)\frac{d\leftsuperscript{B}{H}{D}(\theta)}{d\theta}\Big]^{\vee}
\]
\[
\leftsuperscript{B}{S}{B,D}(\theta) = \Big[\frac{d\leftsuperscript{B}{H}{D}(\theta)}{d\theta}{\leftsuperscript{D}{H}{B}(\theta)}\Big]^{\vee}
\]
\[
\leftsuperscript{D[B]}{S}{B,D}(\theta) = 
\begin{bmatrix}
    \frac{d\leftsuperscript{B}{o}{D}(\theta)}{d\theta}\\
    \big[ \frac{d\leftsuperscript{B}{R}{D}(\theta)}{d\theta}\leftsuperscript{D}{R}{B}(\theta \big]^\vee
\end{bmatrix}^\vee
\]

The relative acceleration between the two links B and D can be obtained by differentiating Equation (\ref{relative_velocity_joint_motion_subspace}):

\begin{equation}
    \leftsuperscript{X}{\dot{\vv}}{B,D} = \iddt{S}\dot{\theta} + S\ddot{\theta} = S\ddot{\theta},
\end{equation}
where $\ddot{\theta}$ is the joint acceleration, and where we assumed that the motion subspaces are independent from the joint configuration: $\frac{d \leftsuperscript{X}{S}{B,D}}{d\theta} = 0$.
 

\subsubsection{Links}\label{sssec_links}
A multi-body system is also composed of $n_L$ links, which describe the mass and inertial properties of a robot through the inertia matrix. A link represents the rigid body that make up the robot's structure, whose motion is constrained by its associated joints. 

\subsubsection{Topology}\label{sssec_topology}
A multi-body system can then be mathematically described as a graph $G = \big( \mathcal{L}, \mathcal{J}, \mathcal{W} \big)$, where $\mathcal{L}$ is the set of all links (the vertices of the graph), $\mathcal{J}$ is the set of joints (the undirected edges of the graph), and $\mathcal{W}$ is the inertial frame called floating base and required by Newton's laws in case the base of the system does not have a fixed, constant position.\\

\begin{defn}[Path]
    A path from a link $B$ to a link $D$ of length $d$ is an ordered sequence of $d$ links $(L_1, L_2, \ldots, L_d)$, such that:

    \begin{equation}\label{def_path}
        L_1 = B,\qquad L_d = D\qquad \forall i\in1,2,\cdots,d \quad \big\{ L_i, L_{i+1} \big\} \in \mathcal{J}
    \end{equation}
    We also define with $\pi_B(L)$ the \textit{unique} set of links that lie on the path connecting $L$ with $B$.
\end{defn}

\begin{defn}[Link frame]
    We assume that each link $L \in \mathcal{L}$ is associated with a link-fixed frame, called \textit{link frame}. Essentially, it is a coordinate system assigned to the link to define its position and orientation in space, such that the description of other connected links becomes immediate with the proper homogeneous transformations.
\end{defn}

\begin{defn}[Parent link]
Assuming that $B \in \mathcal{L}$ is the base link of the model, the parent function
\begin{equation}\label{def_parent_link}
    \lambda_B: \mathcal{L}/B \mapsto \mathcal{L}
\end{equation}
maps each link $L\in \mathcal{L}$ to its parent, with the exclusion of the base link. Since, by definition, the base link $B$ is the root of the tree, it does not have a parent link.
\end{defn}


\subsection{Relative Forward Kinematics}
Let us denote respectively with $s\inrn{n}$ and $\dot{s}\inrn{n}$ the positions and velocities of all joints of the system. It is now possible to define the relationships between the joint positions $s$ of the multi-body system and the relative positions, velocities and accelerations of the bodies composing the system. 
The \textit{relative forward kinematics} is the relative position of two arbitrary bodes $L,D \in \mathcal{L}$, and it is defined as

\begin{equation}\label{forward_kinematics}
    \leftsuperscript{L}{H}{D}(s) = \leftsuperscript{L}{H}{\lambda_D(L)}\leftsuperscript{\lambda_D(L)}{H}{\lambda^2_D(L)}\cdots\leftsuperscript{\lambda^{d-1}_D(L)}{H}{\lambda^d_D(L)}.
\end{equation}
The transform between two nearby links $L$ and $\lambda_D(L)$ is given by the joint model, described in section \ref{sssec_joints}.

As done in \ref{relative_velocity_joint_motion_subspace}, in general the left-trivialised velocity $\leftsuperscript{D}{\vv}{L,D}(s, \dot{s})$ can be written as

\begin{equation}\label{relative_left_trivialised_velocity_jacobian}
    \leftsuperscript{D}{\vv}{L,D}(s, \dot{s}) = \leftsuperscript{D}{S}{L,D}(s)\,\dot{s}(s),
\end{equation}
where $\leftsuperscript{D}{S}{L,D}(s) \inrn{6\times n}$ is the relative left-trivialised \textit{Jacobian}.


\subsection{Positions and velocities}\label{ssec_multibody_lagrangian}
The equations of motion that describe the dynamics of the system are a second-order differential system. This requires the definition of a \textit{state}, which is essentially composed of positions and velocities. 
The configuration of the free-floating system could be described by a set containing all link poses, but the presence of joints that introduce motion constraints allows the description of the system configuration through a pair composed of the pose of a \textit{base link} and the \textit{generalised joint positions}. These two methods of describing the system are respectively knows as \textit{maximal coordinates} and \textit{reduced coordinates}, and the last convention is used throughout the thesis.
Using reduced coordinates, the \textit{generalised position} and \textit{generalised velocity} of the floating-base multi-body system are:

\begin{equation}\label{state_reduced_coordinates}
    \begin{cases}
        q = (\leftsuperscript{W}{H}{B}, s) \in \mathcal{Q} = SE(3) \times \mathbb{R}^n\\
        \dot{q} = \big( \leftsuperscript{W}{\dot{H}}{B}, \dot{s} \big) \in \mathcal{V}
    \end{cases},
\end{equation}
where $s \inrn{n}$ are the joint positions.

Considering the case where the model is only composed of 1-DoF joints, $n$ is the overall number of internal degrees of freedom, which equals the number of joints $n_j$. We can define the system velocity vector $\nu^B \inrn{6+n_j}$ using the left-trivialised velocity representation and the right or mixed velocity representation:

\begin{align}
    \nu^{B/B} = 
    \begin{bmatrix}
        \leftsuperscript{B}{\vv}{A,B}\\
        \dot{s}
    \end{bmatrix} = 
    \begin{bmatrix}
        \leftsuperscript{A}{R}{B}^T \leftsuperscript{A}{\dot{o}}{B}\\
        (\leftsuperscript{A}{R}{B}^T \leftsuperscript{A}{\dot{R}}{B})^\vee\\
        \dot{s}
    \end{bmatrix} \inrn{n+6}
    \\
    \nu^{A/B} = 
    \begin{bmatrix}
        \leftsuperscript{A}{\vv}{A,B}\\
        \dot{s}
    \end{bmatrix} = 
    \begin{bmatrix}
        \leftsuperscript{A}{R}{B}^T \leftsuperscript{A}{\dot{o}}{B}\\
        (\leftsuperscript{A}{R}{B}^T \leftsuperscript{A}{\dot{R}}{B})^\vee\\
        \dot{s}
    \end{bmatrix} \inrn{n+6}
\end{align}

The pose of a link E with respect to the world frame $W$ is
\begin{equation}
    \leftsuperscript{W}{H}{E}(q) = \leftsuperscript{W}{H}{B}\leftsuperscript{B}{H}{E}(s) =
    \begin{bmatrix}
        \leftsuperscript{W}{R}{B} & \leftsuperscript{W}{o}{B}\\
        0_{1\times3} & 1
    \end{bmatrix}\leftsuperscript{B}{H}{E}(s),
\end{equation}
where $\leftsuperscript{B}{H}{E}(s)$ is the relative pose between the base B and the link E, and it is a function of the joint positions $s$.

The velocity of a link E with respect to the world frame $W$ is the sum of the base velocity and the velocity between the base B and link E:

\begin{equation}\label{velocity_link_world_frame}
    \leftsuperscript{E}{\vv}{W,E} + \leftsuperscript{E}{\vv}{W,B} + \leftsuperscript{E}{\vv}{B,E}
\end{equation}

We can rewrite the velocity between the base B and link E as the sum of the velocities between adjacent links in the path $\pi_B(E)$ connecting B to E:

\begin{equation}
    \begin{split}
        \leftsuperscript{E}{v}{W,E} &= \leftsuperscript{E}{v}{W,B} + \sum_{L_i \in \{\pi_B(E)/B\}} \leftsuperscript{E}{\vv}{\lambda(L_i), L_i}\\
        &= \leftsuperscript{E}{\vv}{W,B} + \sum_{L_i \in \{\pi_B(E)/B\}} \leftsuperscript{E}{X}{L_i}\leftsuperscript{L_i}{\vv}{\lambda(L_i), L_i}\\
        &= \leftsuperscript{E}{\vv}{W,B} + \sum_{L_i \in \{\pi_B(E)/B\}} \leftsuperscript{E}{X}{L_i}\leftsuperscript{L_i}{S}{\lambda(L_i), L_i}(s)\dot{s},
    \end{split}
\end{equation}
where the expression for the relative velocity between two adjacent links is given by \ref{relative_left_trivialised_velocity_jacobian}.

In matrix form, the latter equation can be written as

\begin{equation}\label{velocity_link_world_frame_matrix}
    \leftsuperscript{E}{v}{W,E} =
    \begin{bmatrix}
        \leftsuperscript{E}{X}{X} & \leftsuperscript{E}{S}{B,E}(s)
    \end{bmatrix}
    \begin{bmatrix}
        \leftsuperscript{X}{\vv}{W,B}\\
        \dot{s}
    \end{bmatrix},
\end{equation}
where $X$ is a placeholder for the velocity representation used for the system velocity.

By introducing the matrix $\mathbb{S}_{B,E}\inrn{6\times n}$, whose $i^th$ column is defined as

\begin{equation}\label{joint_motion_subspace_matrix}
    \leftsuperscript{E}{S}{B,E}^{(:,i)}(s) =
    \begin{cases}
        \leftsuperscript{E}{X}{L}\leftsuperscript{L}{S}{\lambda(L),L}(s) \qquad\text{if}\;L\in \{ \pi_B(E)/B \}\\
        0_6\qquad\qquad\qquad\text{otherwise}
    \end{cases}
\end{equation}

We can then introduce the floating-base Jacobian of link E as

\begin{equation}\label{floating_base_jacobian}
    \leftsuperscript{E}{J}{W,E/X}(q) = 
    \begin{bmatrix}
        \leftsuperscript{E}{X}{X} & \leftsuperscript{E}{S}{B,E}(s)
    \end{bmatrix} \inrn{6\times (6+n)}
\end{equation}

to rewrite Equation (\ref{velocity_link_world_frame_matrix}) as 

\begin{equation}
    \leftsuperscript{E}{\vv}{W,E} = \leftsuperscript{E}{J}{W,E/X}(q) \leftsuperscript{X}{\nu}{}
\end{equation}

The duality between $6D$ velocities and $6D$ forces also propagates to the definition of the floating-base Jacobian. If we consider a
contact wrench $\leftsuperscript{C_i}{f_i}{}$, which could, for example, be the combination of forces and moments
applied to the frame $C_i = (o_{C_i}, [A])$ associated with contact point i of link E, whose pose
is defined by the transform $\leftsuperscript{E}{H}{C_i}$, we can use the Jacobian defined in Equation (\ref{floating_base_jacobian}) to
project the wrench to the floating-base configuration space as:

\begin{equation}
    \leftsuperscript{C_i}{J}{W,L/X}^T(q) \leftsuperscript{C_i}{f}{} = \big( \leftsuperscript{C_i}{X}{L}\leftsuperscript{L}{J}{W,L/X}(q) \big)^T \leftsuperscript{C_i}{f}{} \inrn{6+n}
\end{equation}
In the particular example where the frame $C_i$ is time-varying, the most appropriate contact Jacobian is $\leftsuperscript{C_i}{J}{W,L}$ 
instead of $\leftsuperscript{C_i}{J}{W,C_i}$, regardless of the system’s velocity representation X.

\subsection{Lagrangian dynamics}

The Lagrangian for a free-floating system can be obtained as the sum of the Lagrangian of each rigid body:

\begin{subequations}
\begin{equation}
    l(q, \nu) = k(q, \nu) - U(q),
\end{equation}    
\begin{equation}
    k(q, \nu) = \sum_{L \in \mathcal{L}} k_L(q, \nu) = \sum_{L \in \mathcal{L}} \leftsuperscript{L}{v}{W,L}^T \mathbb{M}_L \leftsuperscript{L}{v}{W,L} = \frac{1}{2} \nu^T M(s) \nu,
\end{equation}
\begin{equation}
    U(q) =
    \sum_{L \in \mathcal{L}} U_L(q) =
    -\sum_{L \in \mathcal{L}} 
    \begin{bmatrix}
        \leftsuperscript{W}{g}{}^T & 0
    \end{bmatrix}
    \leftsuperscript{W}{H}{L}(q)
    \begin{bmatrix}
        m_L\leftsuperscript{L}{c}{L}\\
        m_L
    \end{bmatrix} = -
    \begin{bmatrix}
        g^T & 0_{1\times3}
    \end{bmatrix} 
    m\leftsuperscript{W}{H}{B}
    \begin{bmatrix}
        \leftsuperscript{B}{c(s)}{}\\
        1
    \end{bmatrix},
\end{equation}
\end{subequations}
where the last equalities are the simplified versions of the kinetic and potential energy, respectively, due to the use of reduced coordinates.

In the expression for the kinetic energy, $M(s) \inrn{(n+6) \times (n+6)}$ is the system's mass matrix, defined as:

\begin{equation}\label{def_mass_matrix}
    M(s) = \sum_{L \in \mathcal{L}} J_L^T(s) \mathbb{M}_L J_L(s),
\end{equation}

$m$ is the total mass of the system:

\begin{equation}
    m \defeq \sum_{L \in \mathcal{L}} m_L,
\end{equation}

and $\leftsuperscript{B}{c}{}(s)$ is the total centre of mass of the system.

The equations of motion for the system are obtained by applying the Euler-Lagrange equations to the Lagrangian $l(q, \nu)$, and are given by:

\begin{equation}\label{def_equations_of_motion_multibody}
    M(s)\dot{\nu} + C(q, \nu)\dot{\nu} + G(q) =
    \begin{bmatrix}
        0_{6\times1}\\
        \tau
    \end{bmatrix}
    + \sum_{L \in \mathcal{L}} J_L^T(s) f_L^{\text{ext}}.
\end{equation}

We recall that:

\begin{equation}
    \begin{aligned}
        &M(q) = \sum_{L \in \mathcal{L}} J_L^T(s) \mathbb{M}_L J_L(s),\\
        &C(q, \nu) = \sum_{L \in \mathcal{L}} J_L^T \Big[ \big( v_L \crossrn {}_L\mathbb{M}_L + {}_L\mathbb{M}_L v_L \times \big) J_L + {}_L\mathbb{M}_L \dot{J_L} \Big], \\
        &G(q) = -M(q)
        \begin{bmatrix}
            \leftsuperscript{W}{R}{B}^T \leftsuperscript{W}{g}{}\\
            0_{3\times1}\\
            0_{n\times1}
        \end{bmatrix},
    \end{aligned}
\end{equation}
where $C(q, \nu)$ is the \textit{Coriolis matrix} and $G(q)$ is the \textit{gravity vector}.
The term $\tau$ is the vector of joint torques, while $f_L^{\text{ext}}$ is an external force applied to the link $L$.



\chapter{Passive dynamic walking}
Passive dynamic walkers are a class of walking robots which exhibit a stable gait down a small inclined slope without the use of any motors. 
These machines provide an elegant illustration of how proper machine design can generate stable and potentially very energy efficient walking.\cite{Tedrake:2004}
In fact, the passive dynamic walking concept is based on the idea that the robot's dynamics can be exploited to generate a stable walking gait without the need for active control.
The concept can then be extended to actuated systems, where the passive dynamics are used to reduce the control effort required to maintain stability.
The energetics in passive walkers consist of a continuous exchange between kinetic and potential energy, and collisions: the energy lost to friction and collisions when th swing foot enters in contact with the ground is compensated by the gradual convertion of potential energy into kinetic energy during the swing phase, as the walker moves down the slope.

\section{History}
The first researches of passive dynamic walkers occurred in the late 1980s with the research of McGeer \cite{T.McGeer:1989}\cite{T.McGeer:1990a}\cite{T.McGeer:1990b}\cite{T.McGeer:1990c}
He started this work with a simple example of a rimless wheel rolling down a slope \cite{T.McGeer:1990c}. During each step, the wheel gains energy from gravity and loses it when colliding with the terrain. 
McGeer also contributed to the idealisation of a synthetic wheel that has a pin joint at the hub, with all legs removed but two consecutive ones, and cut the rim between the two spokes.
Like a passive pendulum the swing leg swings forward, and at the same time the stance leg rolls forward at a constant speed. 
The initial conditions play a huge role in balancing the energy exchange within the system, to produce a stable gait.
As a follow up to the design of these models, McGeer simulated and built a rigid-legged passive walker that attained bidimensional dynamics due to double leg pairs: the outer legs functioned as one pair of legs, while the inner legs functioned as the second pair.
The walker was not completely passive, as it included motors to fold the feet sideways during the swing phase to prevent toe-stubbing.




\chapter{Zero-order optimisation techniques with focus on genetic algorithms}
\section{Introduction to Zero-Order Optimization}

Zero-order optimization methods, also known as derivative-free or black-box optimization techniques, have gained significant attention in engineering applications where objective function derivatives are unavailable, unreliable, or computationally expensive to obtain\cite{Audet2017}. 
In this case, the function might be non-differentiable or noisy, making gradient-based methods unreliable. While first-order methods rely on gradient information or second-order methods that utilise Hessian matrices, zero-order methods only use function evaluations: $f(x)$, where $x$ is the input vector. 
This property makes them suitable for optimizing complex systems such as passive dynamic walkers, where the robot motion's dependency on hardware parameters could be highly nonlinear and difficult to model analytically.\cite{T.McGeer:1990c}\cite{Collins2005}.
Passive dynamic walkers rely on their morphology to achieve stable walking patterns without external inputs on the motors, making their behaviour highly sensitive to parameter configurations\cite{Kuo1999}\cite{Garcia1998}. 
The dynamics of these systems are often complex, where small perturbations in the parameters such as leg length, mass distribution, or foot curvature can lead to highly different walking behaviors, from stable periodic gaits to falling motion\cite{Goswami1998}.

\section{Overview of Zero-Order Optimization Methods}
\subsection{Direct Search Methods}
Direct search methods are a class of zero-order optimization techniques that explore the search space in an iterative way without the need for gradient information. In fact, they represent some of the earliest and most intuitive zero-order optimization approaches. 
For example, the Nelder-Mead method\cite{Nelder1965} remains one of the most used direct search methods due to its simplicity and effectiveness for low-dimensional problems, despite being one of the oldest algorithms on the matter. 
The method maintains a simplex (a generalization of a triangle in n-dimensional space) and iteratively replaces the worst vertex with a better point through reflection, expansion, or contraction operations.

Pattern search methods\cite{Hooke1961}\cite{Torczon1997} systematically explore the parameter space using a predefined pattern of points.
In the past, these methods have been applied to numerous robotics optimization systems, including the work by Chestnutt et al.\cite{Chestnutt2005} for footstep planning in humanoid robots. 
However, their application to passive dynamic walkers has been limited due to the discontinuous nature of the objective function.
Lastly, grid search and random search represent even simpler approaches compared to the prevously mentioned ones, where they systematically or randomly sample the parameter space.
While conceptually simple and intuitive, these methods become computationally limiting when dealing with high-dimensional problems due to the curse of dimensionality.

\subsection{Model-Based Methods}
Model-based optimization methods improve search efficiency by constructing alternative models for the objective function.
One of the earliest approaches by Box and Wilson\cite{Box1951} fits a polynomial model to sampled data points, using this approximation to identify promising regions for further exploration.
Bayesian optimization, for example, takes a probabilistic approach, modeling the objective function with Gaussian processes to balance exploration of uncertain regions with exploitation of known promising areas.
This method has shown to be particularly adequate for optimising passive dynamic walkers.
For example, Calandra\cite{Calandra2016} applied it to gait optimization for a bipedal robot, achieving faster convergence compared to other zero-order methods.

\subsection{Population-Based Methods}
Population-based methods maintain and evolve a collection of candidate solutions.
Some algorithms of this type include:

\begin{itemize}
    \item Particle swarm optimisation (PSO)\cite{Kennedy1995}
    \item Differential evolution\cite{Storn1997}
    \item Covariance Matrix Adaptation Evolution Strategy (CMA-ES)\cite{Hansen2001}
\end{itemize}

\section{Genetic Algorithms: Principles and Applications}
\subsection{Fundamental Concepts of Genetic Algorithms}
Genetic algorithms (GAs), introduced by Holland (1975) and popularized by Goldberg (1989), represent a class of evolutionary computational methods inspired by natural selection.
They operate on a population of candidate solutions (individuals), applying various selection, crossover, and mutation operators to evolve increasingly fit individuals over successive generations.

The basic genetic algorithm follows these key steps:
\begin{enumerate}
    \item \textbf{Initialisation}: Generate a initial population of random candidate solutions, or an initial guess.
    \item \textbf{Evaluation}: Calculate the fitness of all individuals inside the population using the objective function.
    \item \textbf{Selection}: Choose individuals for reproduction based on their fitness: the higher the better.
    \item \textbf{Crossover}: Combine genes from the selected parents to create a new offspring as a new candidate solution.
    \item \textbf{Mutation}: Introduce random changs to the offspring to maintain diversity.
    \item \textbf{Replacement}: Form a new population from the offspring and possibly some parents.
    \item \textbf{Termination}: Stop the algorithm when a termination criterion is met, such as a maximum number of generations or a fitness threshold.
\end{enumerate}

The representation of solutions as `chromosomes` is crucial in genetic algorithms.
For the particular case of the parameter optimisation of a passive walker, floating point representations have largely replaced traditional binary encodings, as they more naturally represent continuous parameters such as link lengths, masses, and joint properties.
Specialized genetic operators for real-valued GAs, such as arithmetic crossover and Gaussian mutation, have demonstrated superior performance for engineering optimization problems\cite{Deb1995}.

\subsection{Advanced GA Techniques Relevant to PDW Optimization}
Several advanced GA techniques have emerged that address challenges particularly relevant to passive dynamic walker optimization:

\begin{itemize}
    \item \textbf{Multi-objective genetic algorithms}, such as NSGA-II\cite{Deb2002}, enable the simultaneous optimization of multiple objectives.
    \item \textbf{Constraint handling techniques} address the many constraints in PDW design. Approaches include penalty functions, repair operators, and specialized constraint-preserving operators.
    \item \textbf{Niching and diversity preservation} techniques help maintain population diversity, which is potentially important for exploring the fitness domains with numerous local optima, which is a common feature of PDW parameter spaces.
    \item \textbf{Hybrid algorithms} combine GAs with local search methods in memetic algorithms.
\end{itemize}



\subsection{GA Applications in Passive Dynamic Walker Optimisation}
The application of genetic algorithms to passive dynamic walker optimisation has evolved considerably over the past three decades.
Some of the earliest works on this matter came from McGeer's group,  using genetic algorithms to optimize a simple compass gait walker for energy efficiency.
A considerable upgrade came with \cite{Collins2001}, who used genetic algorithms to optimize the physical parameters of a three-dimensional, kneed, two-legged passive dynamic walker.
Their work proved that genetic algorithms could efficiently search the complex parameter space of 3D walkers, finding configurations that produced stable stable gaits on slopes without active control.

Parametric optimization of more complex models was undertaken by Trifonov\cite{Trifonov2008}, who used GAs to optimize a kneed walker model with arms.
This work emphasises the ability of genetic algorithms to work well with the increased dimensionality and complexity introduced by additional degrees of freedom on a robot.

\subsection{Advantages of GAs for PDW Optimization}
Genetic algorithms come with numerous advantages that make them particularly suitable for the objectives set by this thesis:

\begin{itemize}
    \item \textbf{Global exploration}: as mentioned earlier, GAs can navigate the complex parameter space of a passive passive walker, with multiple degrees of freedom, as numerous local optima may be present.
    \item \textbf{Parallelizability}: GAs can be efficiently implemented on modern computing architectures, with fitness evaluations of different individuals performed concurrently.
    This derives from a natural independence between individuals in the population.
    \item \textbf{Flexibility in objective function formulation}: it enables the usage of different performance metrics without requiring their derivatives. This allows the optimisation of complex objectives such as mechanical energy efficiency, walking speed and so on.
\end{itemize}


\section{Implementation Considerations for the ptimization}
\subsection{Fitness Function Design}
How I designed the fitness function

\subsection{GA Parameter Tuning}
Which parameters are inserted in the tuning, and their domains.

\subsection{Convergence Criteria}
How the algorithm stops 






\chapter{Stability analysis using Poincaré maps in robotic systems}
\section{Introduction}
Stability analysis is an important aspect of dynamical systems theory, particularly in the study of nonlinear systems, where the primary goal is to determine the stability of equilibrium points and periodic orbits.
Poincaré maps reduce the complexity of continuous dynamical systems by examining their behavior on a lower-dimensional subspace, effectively converting a continuous-time system into a discrete-time system, by using the concept of Poincaré sections. 
This reduction in dimensionality often facilitates the analysis and visualisation of complex behaviours.

A Poincaré map is constructed by first defining a hypersurface $\Sigma$, known as the Poincaré section, which transversely intersects the trajectories of a dynamical system. 
For an $n$-dimensional system, the Poincaré section is typically of dimension $n-1$. 
The map $P: \Sigma \longrightarrow \Sigma$ takes a point $x_0$ on the section and maps it to the point $x_1$, where the trajectory starting from $x_0$ first returns to the section.
In a robotic environment, for example, typical Poincaré sections are chosen to match with specific events in the robot's operation cycle, such as:

\begin{itemize}
    \item The instant of foot impact with the ground in walking robots
    \item The instant of interaction between a chosen end-effector and an object during manipulation tasks
    \item State transitions in hybrid robotic systems, and so on.
\end{itemize}


Formally, for a robotic system described by the differential equation $\dot{x} = f(x, u)$,
where $x \inrn{n}$ represents the state of the robot and $u \inrn{m}$ represents control inputs, the Poincaré map is defined as:
\[
P(x) = \phi(\tau(x), x,u(\cdot)),
\]
where $\phi(\tau(x), x, u(\cdot))$ represents the solution at time $t$ with initial condition $x$ and control input trajectory $u(\cdot)$, and $\tau(x)$ is the first return time (for example the smallest $t > 0$ such that $\phi(\tau(x), x, u(\cdot)) \in \Sigma$).

In robotics applications, fixed points often represent:
\begin{itemize}
    \item Stable walking gaits in legged robots
    \item Cyclic manipulation tasks
    \item Rhythmic movements in bio-inspired robots
    \item Stable flight patterns in flapping-wing aerial robots
\end{itemize}


\section{Stability Analysis via Poincaré Maps for Robots}
The stability of fixed points in the Poincaré map directly relates to the stability of the corresponding periodic orbits in the continuous robotic system. The local stability is determined by examining the eigenvalues of the Jacobian matrix $DP(x^{*})$ evaluated at the fixed point.
The stability criteria for fixed points are:
\begin{itemize}
    \item If $|\lambda_i| < 1 \forall i$, the fixed point is a stable point of equilibrium, indicating that the robot will return to the periodic motion when applying small perturbations to its state.
    \item If $|\lambda_i| > 1$ for at least one $i$, the fixed point is unstable. This means that the robot will deviate from the desired motion pattern when perturbed.
    \item If all $|\lambda_i| = 1$, higher-order analysis is needed to determine stability, but this typically results in an unstable fixed point.
\end{itemize}

